\section{Domain Task 2: Phân tích Tỷ lệ lấp đầy theo thời gian}

Câu hỏi nghiệp vụ: ``Tỷ lệ lấp đầy phòng thay đổi như thế nào qua các
tháng, và tính mùa vụ ảnh hưởng thế nào đến từng quận?''

Mục đích: Xác định xu hướng mùa cao điểm, thấp điểm để từ đó có cái nhìn
chi tiết về hiệu suất kinh doanh qua thời gian của từng khu vực.

\subsection{Biểu đồ 2: Số lượng reviews trung bình theo tháng của từng quận (Multi-line Chart)}

\subsubsection*{A. Thiết kế Idiom}

\begin{longtable}{@{} p{2.8cm} p{12cm} @{}}
\toprule
\textbf{Đặc điểm} & \textbf{Chi tiết} \\
\midrule
\endhead

\bottomrule
\endfoot

\textbf{Idiom} & Multi-line Chart (Biểu đồ đa đường) \\
\addlinespace[0.2cm]

\textbf{What} & Tháng (Month): Ordinal (O), Quận (Neighbourhood Group Cleansed): Categorical (C), Số lượng review trung bình (AvgReviewOfMonth): Quantitative (Q) \\
\addlinespace[0.2cm]

\textbf{Why} & produce (derive) $\rightarrow$ browse, lookup $\rightarrow$ summarize, compare
\begin{itemize}[nosep, leftmargin=*, topsep=0.1cm]
    \item Tạo ra (derive) biến số lượng review trung bình theo tháng để làm thước đo đại diện cho lưu lượng khách lưu trú thực tế.
    \item Duyệt (browse) dọc theo trục thời gian để khám phá xu hướng để lại đánh giá của khách hàng qua các mùa, đồng thời tra cứu (lookup) chính xác lượng review tại các ``đỉnh'' (peak).
    \item Tóm tắt (summarize) mùa cao điểm/thấp điểm và so sánh (compare) mức độ sôi động về mặt tương tác của khách thuê giữa lõi trung tâm và các quận vùng ven.
\end{itemize} \\
\addlinespace[0.2cm]

\textbf{How} & \textbf{Encode:}
\begin{itemize}[nosep, leftmargin=*, topsep=0.1cm]
    \item Mark: Line (Đường nối các điểm)
    \item Channel:
    \begin{itemize}[nosep, leftmargin=0.5cm]
        \item PosX: Tiến trình thời gian (Tháng 1 -- 12)
        \item PosY: Độ lớn của số lượng review (AvgReviewOfMonth)
        \item Color Hue: Phân biệt 5 quận (Categorical Blue-Orange-Red-Teal-Green)
    \end{itemize}
\end{itemize}

\vspace{0.2cm}
\textbf{Manipulate:}
\begin{itemize}[nosep, leftmargin=*, topsep=0.1cm]
    \item Hover + Tooltip: Di chuột vào các điểm gãy để đọc con số review trung bình chính xác.
    \item Highlight: Rê chuột vào một đường trên Legend để làm mờ các đường khác, giúp theo dõi riêng biệt một quận.
\end{itemize}

\vspace{0.2cm}
\textbf{Reduce:}
\begin{itemize}[nosep, leftmargin=*, topsep=0.1cm]
    \item Filter: Kết hợp bộ lọc Tháng/Năm (Month/Year) để xem xu hướng review của các năm cụ thể. Có thể lọc ra các tháng quá gần ngày cào dữ liệu mà chưa có review để đảm bảo khách quan.
\end{itemize} \\
\addlinespace[0.2cm]

\textbf{Scale} &
\begin{itemize}[nosep, leftmargin=*, topsep=0pt]
    \item Main key: 12 (Tháng)
    \item Categorical key: 5 (Nhóm quận)
    \item Quantitative range: Từ 0 đến khoảng 2400 (reviews)
\end{itemize} \\
\end{longtable}

\begin{figure}[H]
    \centering
    \safeincludegraphics[width=0.9\linewidth]{charts/domain_task2_chart2.png}
    \caption{Hình 2.2. Biểu đồ số lượng reviews trung bình theo tháng của từng quận}
    \label{fig:domain-task2-chart2}
\end{figure}

\subsubsection*{B. Đánh giá biểu đồ}

\textbf{1. Nguyên lý biểu đạt (Expressiveness):}

\begin{itemize}
    \item Biểu đồ tuân thủ hoàn toàn bản chất của dữ liệu chuỗi thời gian phân nhóm.
    \item Đánh giá mức độ quan trọng và phân bổ kênh theo Mackinlay:
    \begin{itemize}
        \item Ưu tiên 1 -- Vị trí trục tung (PosY): Số lượng review trung bình (Q) là mục tiêu quan sát cốt lõi để đo lường mức độ hoạt động. Nó được phân bổ một cách hoàn hảo cho kênh Vị trí (Position) --- kênh có sức mạnh biểu đạt số liệu định lượng chính xác nhất theo quy luật thị giác.
        \item Ưu tiên 2 -- Vị trí trục hoành (PosX): Tháng (O) tạo nên tính tuần tự của thời gian. Kênh PosX giúp biểu diễn mũi tên thời gian từ trái sang phải, thuận với quy luật đọc tự nhiên của mắt.
        \item Ưu tiên 3 -- Màu sắc (Color Hue): Thuộc tính Quận (C) được dùng để so sánh chéo. Việc gán sắc độ màu (Hue) cho dữ liệu phân loại là sự lựa chọn tối ưu thứ hai của Mackinlay khi kênh Vị trí đã được sử dụng.
    \end{itemize}
\end{itemize}

\textbf{2. Nguyên lý hiệu quả (Effectiveness):}

\begin{itemize}
    \item Độ chính xác (Accuracy): Do PosY là một thang đo chung bắt đầu từ 0 và phát triển theo chiều dài tuyến tính, định luật Stevens đánh giá kênh này có độ lỗi Log\_error cực thấp ($n = 1.0$). Mắt người có thể dễ dàng so sánh khoảng cách giữa đường màu cam (Brooklyn) và màu xanh lá (Staten Island) và nhận thức được tỷ lệ chênh lệch khổng lồ (gấp hàng chục lần).
    \item Khả năng phân biệt (Discriminability): Palette màu gồm 5 màu phân biệt rõ ràng, không gây quá tải cho bộ nhớ ngắn hạn của mắt (vì $< 7$ màu). Hiện tượng cắt nhau (line crossing) có xuất hiện nhẹ ở giai đoạn đầu năm giữa Manhattan (đỏ) và Brooklyn (cam), nhưng không gây ra tình trạng Clutter (rối mắt) vì quỹ đạo của các đường giãn cách rất rộng ở phần lớn thời gian còn lại.
    \item Khả năng tách biệt (Separability): Kênh không gian (Pos) và kênh màu sắc (Color) tách biệt hoàn toàn. Độ cao thấp của điểm dữ liệu không làm ảnh hưởng đến khả năng nhận diện màu sắc của mắt.
\end{itemize}

\subsubsection*{C. Phân tích biểu đồ (Insight)}

\begin{itemize}
    \item Sự chênh lệch về mức độ tương tác: Brooklyn (cam) và Manhattan (đỏ) thống trị hoàn toàn về lượng khách lưu trú và để lại review, vượt xa 3 quận còn lại. Queens (xanh mòng két) nằm ở mức trung bình, trong khi Bronx và Staten Island gần như ``chạm đáy'' trục hoành với lượng review lẹt đẹt suốt cả năm.
    \item Tính mùa vụ thể hiện qua Review: Có một xu hướng tăng trưởng cực kỳ rõ rệt bắt đầu từ mùa xuân (tháng 3-4) và đạt đỉnh điểm vào cuối mùa hè / đầu mùa thu (tháng 8 -- tháng 10). Lượng review của Brooklyn tạo đỉnh chóp nón lên tới hơn 2200 đánh giá/tháng vào tháng 9. Sau tháng 10, lượng tương tác giảm mạnh do bước vào mùa đông.
\end{itemize}
